%===============================================================================
% LLMWiki Technical Specification — Volume 08: API Reference
%===============================================================================
\chapter{API Reference}
\label{cha:api}

This volume specifies the public API contracts: REST endpoints, gRPC services, WebSocket protocol, authentication, rate limiting, and SDKs. All APIs are versioned; breaking changes only in major versions.

\section{API Overview}
\label{sec:api:overview}

\begin{itemize}
  \item \textbf{Base URL}: \texttt{http://<host>:8080/api/v1}
  \item \textbf{Protocol}: HTTP/1.1, h2 (gRPC), WebSocket
  \item \textbf{Format}: JSON (request/response), SSE (streaming)
  \item \textbf{Auth}: Bearer token (OIDC), API Key, or none (dev)
  \item \textbf{Versioning}: URL path \texttt{/v1/}, header \texttt{Accept: application/vnd.llmwiki.v1+json}
\end{itemize}

\section{Authentication}
\label{sec:api:auth}

\subsection{Bearer Token (OIDC)}
\label{sec:api:oidc}

\begin{verbatim}
Authorization: Bearer <access_token>
\end{verbatim}

Token validated against configured OIDC issuer. Claims mapped to \texttt{UserContext}:
\begin{verbatim}
{
  "sub": "user_123",
  "email": "user@example.com",
  "roles": ["admin", "analyst"],
  "org_id": "org_456",
  "permissions": ["query", "ingest", "admin"]
}
\end{verbatim}

\subsection{API Key}
\label{sec:api:apikey}

\begin{verbatim}
Authorization: ApiKey <key_id>.<secret>
# or header
X-API-Key: <key_id>.<secret>
\end{verbatim}

API keys managed via \texttt{/admin/api-keys} (admin only). Scoped permissions.

\subsection{Anonymous (Development)}
\label{sec:api:anon}

\texttt{auth: "none"} in config disables auth. Only for local development.

\section{Core Endpoints}
\label{sec:api:endpoints}

\subsection{Health}
\label{sec:api:health}

\begin{verbatim}
GET /health/live
GET /health/ready
GET /health/deep

Response (ready):
{
  "status": "ready",
  "checks": {
    "graph": "ok",
    "index": "ok",
    "evidence_gate": "ok",
    "wal": "ok"
  },
  "version": "0.1.0",
  "repo_id": "prod-main"
}
\end{verbatim}

\subsection{Query}
\label{sec:api:query}

\paragraph{Execute Query (Non-streaming)}
\begin{verbatim}
POST /query
Content-Type: application/json

{
  "query": "How does the payment service handle retries?",
  "context_budget": 8000,
  "evidence_threshold": 0.72,
  "max_steps": 5,
  "stream": false,
  "plan_options": {
    "max_depth": 3,
    "include_graph": true
  }
}

Response:
{
  "query_id": "q_abc123",
  "answer": "The payment service uses exponential backoff...",
  "citations": [
    {
      "unit_id": "sha256:a1b2c3...",
      "relevance": 0.94,
      "span": {"start_line": 45, "end_line": 67},
      "preview": "fn handle_retry(...) { ... }"
    }
  ],
  "evidence_count": 12,
  "rejected_count": 35,
  "plan": {...},
  "usage": {
    "tokens": 2450,
    "latency_ms": 312,
    "llm_calls": 2
  }
}
\end{verbatim}

\paragraph{Execute Query (Streaming)}
\begin{verbatim}
POST /query/stream
Content-Type: application/json
Accept: text/event-stream

{ "query": "...", "stream": true, ... }

SSE Events:
event: plan_start
data: {"plan_id": "p_123", "steps": 5}

event: step_start
data: {"step_id": "s_1", "kind": "Search", "tool": "hybrid_search"}

event: evidence
data: {"unit_id": "sha256:...", "relevance": 0.94, "preview": "..."}

event: token
data: {"token": "The", "step_id": "s_3"}

event: token
data: {"token": " payment", "step_id": "s_3"}

event: step_complete
data: {"step_id": "s_1", "output": {"results": [...]}}

event: plan_complete
data: {"answer": "...", "citations": [...], "usage": {...}}
\end{verbatim}

\subsection{Ingestion}
\label{sec:api:ingest}

\paragraph{Ingest Path}
\begin{verbatim}
POST /ingest
Content-Type: application/json

{
  "path": "/data/source-code",
  "mode": "incremental",  # or "full"
  "parsers": ["python", "markdown", "yaml"],
  "options": {
    "chunk_size": 512,
    "overlap": 50,
    "embed": true
  }
}

Response:
{
  "job_id": "job_xyz789",
  "status": "accepted",
  "estimated_units": 5000
}
\end{verbatim}

\paragraph{Ingest Status}
\begin{verbatim}
GET /ingest/{job_id}

Response:
{
  "job_id": "job_xyz789",
  "status": "running",  # pending|running|completed|failed
  "progress": 0.65,
  "units_processed": 3250,
  "units_total": 5000,
  "added": 3200,
  "modified": 45,
  "removed": 5,
  "errors": 0,
  "started_at": "2024-01-15T10:30:00Z",
  "updated_at": "2024-01-15T10:31:45Z"
}
\end{verbatim}

\paragraph{Ingest History}
\begin{verbatim}
GET /ingest/history?limit=20
\end{verbatim}

\subsection{Repository Management}
\label{sec:api:repo}

\begin{verbatim}
GET /repo              # Current repo info
POST /repo/init        # Initialize new repo
POST /repo/clone       # Clone from git/remote
GET /repo/stats        # Unit/graph/index counts
POST /repo/verify      # Integrity check
POST /repo/compact     # Reclaim space, rebuild indices
POST /repo/backup      # Trigger backup
GET /repo/wal          # WAL entries (paginated)
\end{verbatim}

\subsection{Graph Operations}
\label{sec:api:graph}

\paragraph{Traverse}
\begin{verbatim}
POST /graph/traverse
{
  "start": "sha256:...",
  "pattern": {
    "edge_types": ["CALLS", "REFERENCES"],
    "node_types": ["Function", "Class"],
    "direction": "out",
    "max_hops": 3,
    "limit": 100
  }
}

Response:
{
  "nodes": [...],
  "edges": [...],
  "subgraph_id": "sg_123"
}
\end{verbatim}

\paragraph{Entity Lookup}
\begin{verbatim}
GET /graph/entity/{entity_id}
GET /graph/entity/by-name?name=process_order&type=Function
\end{verbatim}

\paragraph{Neighborhood}
\begin{verbatim}
GET /graph/neighborhood?ids=sha256:a,sha256:b&hops=2
\end{verbatim}

\subsection{Index Operations}
\label{sec:api:index}

\begin{verbatim}
POST /index/rebuild        # Full rebuild
POST /index/reembed        # Re-embed (model change)
GET /index/stats           # Vector count, dimension, memory
GET /index/health          # HNSW/BM25 status
\end{verbatim}

\subsection{Unit Access}
\label{sec:api:unit}

\begin{verbatim}
GET /units/{unit_id}           # Full unit
GET /units/{unit_id}/raw       # Raw content
GET /units/batch?ids=...       # Multiple units
GET /units/search?q=...        # Search units (alias for /query)
\end{verbatim}

\subsection{Evidence Gate}
\label{sec:api:evidence}

\begin{verbatim}
POST /evidence/verify
{
  "query": "How does X work?",
  "candidates": ["sha256:a", "sha256:b", ...],
  "threshold": 0.72
}

Response:
{
  "verified": [
    {"unit_id": "sha256:a", "score": 0.94}
  ],
  "rejected": [
    {"unit_id": "sha256:b", "score": 0.45}
  ],
  "calibration": "isotonic"
}
\end{verbatim}

\subsection{Administration}
\label{sec:api:admin}

\begin{verbatim}
GET  /admin/config              # Current config (sanitized)
POST /admin/config/reload       # Reload from file
GET  /admin/metrics             # Prometheus metrics
GET  /admin/pprof/...           # Go-style pprof (if enabled)

# API Key management
POST   /admin/api-keys            # Create
GET    /admin/api-keys            # List
DELETE /admin/api-keys/{key_id}   # Revoke

# Parser management
GET    /admin/parsers             # List loaded parsers
POST   /admin/parsers/reload      # Reload parser configs
\end{verbatim}

\section{WebSocket API}
\label{sec:api:ws}

For real-time, bidirectional communication (CLI, IDE extensions).

\begin{verbatim}
WS /ws/v1

# Client -> Server
{"type": "query", "payload": {"query": "...", "stream": true}}
{"type": "ingest", "payload": {"path": "..."}}
{"type": "subscribe", "payload": {"events": ["ingest.progress", "query.token"]}}

# Server -> Client
{"type": "event", "event": "query.token", "data": {"token": "..."}}
{"type": "event", "event": "ingest.progress", "data": {"progress": 0.5}}
{"type": "result", "request_id": "...", "data": {...}}
{"type": "error", "request_id": "...", "code": "...", "message": "..."}
\end{verbatim}

\section{gRPC API}
\label{sec:api:grpc}

Defined in \texttt{proto/llmwiki/v1/api.proto}. Mirrors REST with strongly-typed messages.

\begin{verbatim}
service LLMWiki {
  rpc Query(QueryRequest) returns (QueryResponse);
  rpc QueryStream(QueryRequest) returns (stream QueryStreamEvent);
  rpc Ingest(IngestRequest) returns (IngestResponse);
  rpc IngestProgress(IngestProgressRequest) returns (stream IngestProgress);
  rpc GraphTraverse(GraphTraverseRequest) returns (GraphTraverseResponse);
  rpc UnitGet(UnitGetRequest) returns (Unit);
  rpc HealthCheck(HealthRequest) returns (HealthResponse);
}
\end{verbatim}

gRPC gateway (transcoding) exposes same REST endpoints automatically.

\section{Error Handling}
\label{sec:api:errors}

\begin{verbatim}
# Standard error response
{
  "error": {
    "code": "EVIDENCE_INSUFFICIENT",
    "message": "No evidence units passed verification threshold",
    "details": {
      "query_id": "q_abc",
      "candidates": 42,
      "verified": 0,
      "threshold": 0.72
    }
  },
  "request_id": "req_xyz",
  "timestamp": "2024-01-15T10:30:00Z"
}
\end{verbatim}

\begin{table}[htbp]
\centering
\caption{Error Codes}
\label{tab:api:errors}
\begin{tabular}{ll}
\toprule
\textbf{Code} & \textbf{Meaning} \\
\midrule
INVALID_REQUEST & Request validation failed \\
UNAUTHENTICATED & Missing/invalid auth \\
FORBIDDEN & Insufficient permissions \\
NOT_FOUND & Resource not found \\
CONFLICT & Resource conflict (e.g., repo exists) \\
EVIDENCE_INSUFFICIENT & Gate rejected all candidates \\
BUDGET_EXCEEDED & Token/time/cost budget exceeded \\
INTERNAL_ERROR & Unexpected server error \\
SERVICE_UNAVAILABLE & Dependency down (embedder, LLM) \\
RATE_LIMITED & Too many requests \\
\bottomrule
\end{tabular}
\end{table}

\section{Rate Limiting}
\label{sec:api:ratelimit}

\begin{verbatim}
Headers on all responses:
X-RateLimit-Limit: 100
X-RateLimit-Remaining: 95
X-RateLimit-Reset: 1705312800

On 429:
Retry-After: 60
\end{verbatim}

Limits configurable per auth type (API key, user, IP, global).

\section{SDKs}
\label{sec:api:sdks}

\begin{table}[htbp]
\centering
\caption{Official SDKs}
\label{tab:api:sdks}
\begin{tabular}{lll}
\toprule
\textbf{Language} & \textbf{Package} & \textbf{Status} \\
\midrule
Python & \texttt{pip install llmwiki} & v0.1 (alpha) \\
TypeScript/JS & \texttt{npm install @llmwiki/sdk} & v0.1 (alpha) \\
Rust & \texttt{cargo add llmwiki-client} & v0.1 (alpha) \\
Go & \texttt{go get github.com/llmwiki/go-sdk} & Planned \\
\bottomrule
\end{tabular}
\end{table}

\paragraph{Python SDK Example}
\begin{verbatim}
from llmwiki import LLMWikiClient

client = LLMWikiClient(
    base_url="http://localhost:8080",
    api_key="key_id.secret"
)

# Simple query
result = client.query("How does retry logic work?")
print(result.answer)
for cite in result.citations:
    print(f"  [{cite.unit_id[:8]}] {cite.preview}")

# Streaming
for event in client.query_stream("Explain the auth flow"):
    if event.type == "token":
        print(event.token, end="", flush=True)
    elif event.type == "plan_complete":
        print("\n---")
        print(f"Evidence: {event.data.evidence_count}")
\end{verbatim}

\section{OpenAPI Specification}
\label{sec:api:openapi}

Full OpenAPI 3.1 spec at \texttt{/openapi.json} and \texttt{/openapi.yaml}.
Generated from source-of-truth \texttt{api/src/routes.rs} (Axum) or \texttt{api/src/main.rs} (Actix).

\section{API Versioning Policy}
\label{sec:api:versioning}

\begin{itemize}
  \item \textbf{v1}: Current stable. No breaking changes.
  \item \textbf{Deprecation}: 6 months notice via \texttt{Sunset} header and changelog.
  \item \textbf{New versions}: \texttt{/v2/} path; both served simultaneously during transition.
  \item \textbf{Experimental}: \texttt{/v1/experimental/} prefix; no stability guarantee.
\end{itemize}

\section{Patent-Relevant Novelty}
\label{sec:api:patent}

\begin{patentclaim}
\textbf{Claim 11:} An API for a semantic operating layer comprising: (a) a query endpoint that returns answers with mandatory evidence citations verified by a cross-encoder gate, (b) an ingestion endpoint that accepts filesystem paths and returns content-addressed unit diffs, (c) a graph traversal endpoint exposing the knowledge graph as a queryable API, (d) a streaming protocol that emits fine-grained events (plan, evidence, token, citation) for real-time UI, and (e) an evidence verification endpoint that exposes the gate as a first-class callable service.
\end{patentclaim}

\section{Summary}
\label{sec:api:summary}

The LLMWiki API exposes the full power of the semantic operating layer through:
\begin{enumerate}
  \item \textbf{REST}: Simple, cacheable, widely compatible
  \item \textbf{gRPC}: High-performance, typed, streaming-native
  \item \textbf{WebSocket}: Bidirectional, real-time, IDE-friendly
  \item \textbf{SDKs}: Idiomatic clients for major languages
\end{enumerate}
All endpoints respect the architectural invariants: evidence closure, content addressing, deterministic replay.

The next volume (Vol.~09) guides developers extending LLMWiki.

%===============================================================================
% End of Volume 08
%===============================================================================
